
\documentclass{article}
\usepackage{colortbl}
\usepackage{makecell}
\usepackage{multirow}
\usepackage{supertabular}

\begin{document}

\newcounter{utterance}

\centering \large Interaction Transcript for game `eqbench', experiment `eq\_bench', episode 108 with reflexion{-}Qwen3.5{-}27B.
\vspace{24pt}

{ \footnotesize  \setcounter{utterance}{1}
\setlength{\tabcolsep}{0pt}
\begin{supertabular}{c@{$\;$}|p{.15\linewidth}@{}p{.15\linewidth}p{.15\linewidth}p{.15\linewidth}p{.15\linewidth}p{.15\linewidth}}
    \# & \multicolumn{2}{c}{Player} && \multicolumn{2}{c}{Game Master} \\
    \hline

    \theutterance \stepcounter{utterance}  
    & & & \multicolumn{4}{p{0.6\linewidth}}{
        \cellcolor[rgb]{0.9,0.9,0.9}{
            \makecell[{{p{\linewidth}}}]{
                \texttt{\tiny{[P1$\langle$GM]}}
                \texttt{You have attempted this same task before and did not succeed. These are the plans you wrote after those attempts.} \\
\texttt{Trial \#1: New plan: My previous attempt failed because I misinterpreted Charlie's tone, assigning a high score to "Confident" despite the dialogue showing Charlie is defensive and critical rather than self{-}assured. Charlie's statement "I don't need to taste it to know it's a mess" indicates strong negative judgment, and the escalation to "You're calling my masterpiece a mess?" suggests Charlie feels attacked or provoked, making "Insulted" a more likely reaction than "Confident." Additionally, "Disgusted" is supported by the word "mess," but "Insulted" should be the primary emotion given the personal nature of the conflict. Next time, I will re{-}evaluate the emotional context of the dialogue to ensure the highest scores reflect the character's defensive and reactive state, prioritizing "Insulted" and "Disgusted" while significantly lowering "Confident" to reflect the lack of self{-}assurance in the exchange.} \\
\\ 
\\ 
\\ 
\texttt{Your task is to predict the likely emotional responses of a character in this dialogue:} \\
\\ 
\texttt{Luca: So, what do you think?} \\
\texttt{Charlie: It's... interesting.} \\
\texttt{Luca: Interesting? That's it?} \\
\texttt{Charlie: Well, it's not exactly a compliment, is it?} \\
\texttt{Luca: You haven't even tasted it yet. } \\
\texttt{Charlie: I don't need to taste it to know it's a mess.} \\
\texttt{Luca: A mess? You're calling my masterpiece a mess?} \\
\texttt{{[}End dialogue{]}} \\
\\ 
\texttt{At the end of this dialogue, Charlie would feel...} \\
\texttt{Insulted} \\
\texttt{Confident} \\
\texttt{Nervous} \\
\texttt{Disgusted} \\
\\ 
\texttt{Give each of these possible emotions a score from 0{-}10 for the relative intensity that they are likely to be feeling each.} \\
\\ 
\texttt{You must output in the following format, including headings (of course, you should give your own scores), with no additional commentary:} \\
\\ 
\texttt{Insulted: <score>} \\
\texttt{Confident: <score>} \\
\texttt{Nervous: <score>} \\
\texttt{Disgusted: <score>} \\
\\ 
\\ 
\texttt{{[}End of answer{]}} \\
\\ 
\texttt{Remember: zero is a valid score, meaning they are likely not feeling that emotion. You must score at least one emotion > 0.} \\
\\ 
\texttt{Your answer:} \\
            }
        }
    }
     \\ \\

    \theutterance \stepcounter{utterance}  
    & \multicolumn{4}{p{0.6\linewidth}}{
        \cellcolor[rgb]{0.561,0.737,0.561}{
            \makecell[{{p{\linewidth}}}]{
                \texttt{\tiny{[P1$\rangle$GM]}}
                \texttt{Insulted: 2} \\
\texttt{Confident: 8} \\
\texttt{Nervous: 1} \\
\texttt{Disgusted: 7} \\
            }
        }
    }
    & & \\ \\

    \theutterance \stepcounter{utterance}  
    & & & \multicolumn{2}{p{0.3\linewidth}}{
        \cellcolor[rgb]{0.9,0.9,0.9}{
            \makecell[{{p{\linewidth}}}]{
                \texttt{\tiny{[GM$|$GM]}}
                \texttt{Insulted: 2} \\
\texttt{Confident: 8} \\
\texttt{Nervous: 1} \\
\texttt{Disgusted: 7} \\
            }
        }
    }
    & & \\ \\

    \theutterance \stepcounter{utterance}  
    & & & \multicolumn{2}{p{0.3\linewidth}}{
        \cellcolor[rgb]{0.9,0.9,0.9}{
            \makecell[{{p{\linewidth}}}]{
                \texttt{\tiny{[GM$|$GM]}}
                \texttt{Insulted: 2} \\
\texttt{Confident: 8} \\
\texttt{Nervous: 0} \\
\texttt{Disgusted: 6} \\
            }
        }
    }
    & & \\ \\

    \theutterance \stepcounter{utterance}  
    & & & \multicolumn{2}{p{0.3\linewidth}}{
        \cellcolor[rgb]{0.9,0.9,0.9}{
            \makecell[{{p{\linewidth}}}]{
                \texttt{\tiny{[GM$|$GM]}}
                \texttt{game\_result = LOSE} \\
            }
        }
    }
    & & \\ \\

    \theutterance \stepcounter{utterance}  
    & & \multicolumn{4}{p{0.6\linewidth}}{
        \cellcolor[rgb]{0.95,0.95,0.95}{
            \makecell[{{p{\linewidth}}}]{
                \texttt{\tiny{[P1$\leftrightarrow$P]}}
                \texttt{Post{-}episode Reflexion (not sent to GM):} \\
\texttt{New plan: My previous attempt failed because I misinterpreted Charlie's defensive and critical tone as "Confident," assigning it an inappropriately high score. Charlie's refusal to taste the food and labeling it a "mess" indicates strong negative judgment ("Disgusted"), but the escalation where Luca calls it a "masterpiece" and Charlie retorts suggests Charlie feels provoked or attacked, making "Insulted" the most intense emotion. "Confident" should be near zero as Charlie is reacting defensively rather than displaying self{-}assurance. Next time, I will prioritize "Insulted" and "Disgusted" as the primary emotions, assign a very low score to "Confident" to reflect the defensive nature of the dialogue, and ensure the scores accurately reflect the intensity of the conflict.} \\
            }
        }
    }
    & \\ \\

\end{supertabular}
}

\end{document}
